\section{Tests Beyond the Benchmark}
\label{sec:beyond}
% SWEVO expansion: generated supplementary tables placed in Sections 8 and 9 (\SuppTable{<label>}). The capture is global;
% the figure of mpce_supp_diag.tex (fig:D-psodyn) is discarded here, so only its tables are stored.
\begingroup\RenewEnviron{figure}[1][]{}\RenewEnviron{figure*}[1][]{}%
\CaptureSuppTables{analysis/mpce_supp_diag.tex}\endgroup
\CaptureSuppTables{analysis/mpce_supp_direction.tex}
% Captions of the direction-resolution tables refer to main-text tables as M-<label> (supplement build); in the main
% text these names are aliases of the labels themselves (resolved from the .aux of the previous run).
\makeatletter
\def\sw@mainalias#1{\@ifundefined{r@M-#1}{\@ifundefined{r@#1}{}{%
  \global\expandafter\let\csname r@M-#1\expandafter\endcsname\csname r@#1\endcsname}}{}}
\sw@mainalias{tab:friedman68}\sw@mainalias{tab:hr-site}
\makeatother
\subsection{Budget Scaling}
\label{sec:budget}
%% generated by mpce_results.py -- do not edit by hand
\begin{table}[!t]
\centering
\caption{Six Largest Benchmark Cases: Mean Wake Loss (\%) and Feasible Runs (\%) at 6,030 Calls with Random and Feasibility-Preserving Initialization, and Average Rank at 6,030, 30,030 and 120,030 Calls (Random Initialization; 30 Seeds)}
\label{tab:feasbudget}
\scriptsize\setlength{\tabcolsep}{2.4pt}
\begin{tabular}{lccccccc}
\toprule
& \multicolumn{2}{c}{Random init.} & \multicolumn{2}{c}{Feasible init.} & \multicolumn{3}{c}{Avg.\ rank} \\
\cmidrule(lr){2-3}\cmidrule(lr){4-5}\cmidrule(lr){6-8}
Method & Loss & Feas. & Loss & Feas. & 6,030 & 30,030 & 120,030 \\
\midrule
PSO-VNS & 4.084 & 100.0 & 3.938 & 100.0 & 1.50 & 1.17 & 2.17 \\
SSA-VNS & 4.678 & 95.0 & \TBD{} & \TBD{} & 4.33 & 4.17 & 4.17 \\
LX-SSA-VNS & 4.802 & 98.9 & \TBD{} & \TBD{} & 5.17 & 5.67 & 5.83 \\
RS-VNS & 4.830 & 93.3 & \TBD{} & \TBD{} & 4.67 & 6.33 & 4.67 \\
VNS & 4.950 & 100.0 & \TBD{} & \TBD{} & 6.50 & 4.83 & 3.17 \\
PSO & 4.181 & 100.0 & \TBD{} & \TBD{} & 1.83 & 2.50 & 5.83 \\
SSA & 4.828 & 80.0 & \TBD{} & \TBD{} & 8.33 & 8.67 & 8.83 \\
LX-SSA & 5.379 & 77.8 & \TBD{} & \TBD{} & 8.67 & 9.83 & 10.00 \\
DE & -- & 1.1 & \TBD{} & \TBD{} & 10.00 & 5.17 & 3.33 \\
MS-SLSQP & 4.522 & 99.4 & \TBD{} & \TBD{} & 4.00 & 6.67 & 7.00 \\
\bottomrule
\multicolumn{8}{l}{\TBD{pending: feas}}
\end{tabular}
\end{table}


On the six largest benchmark cases (random starts), \NBudgetPhrase{} (6,030, 30,030 and 120,030 evaluations: \NBudRankPSOVNSSixK, \NBudRankPSOVNSThirtyK{} and \NBudRankPSOVNSOneTwentyK; Table~\ref{tab:feasbudget}; Fig.~\ref{fig:budget}) and ranks first in \NBudFirstPSOVNSSixK, \NBudFirstPSOVNSThirtyK{} and \NBudFirstPSOVNSOneTwentyK{} of the six cases (exploratory; rank differences not tested); PSO's mean wake loss minus that of PSO-VNS is \NBudGainPSOSixK, \NBudGainPSOThirtyK{} and \NBudGainPSOOneTwentyK~pp (Table~\ref{S-tab:budget-cases}). The order of the other methods depends on the budget: from 6,030 to 120,030 evaluations, PSO falls back from an average rank of \NBudRankPSOSixK{} to \NBudRankPSOOneTwentyK, behind SSA-VNS (\NBudRankSSAVNSOneTwentyK), and VNS moves up from \NBudRankVNSSixK{} to \NBudRankVNSOneTwentyK; at 120,030, \NBudCloseListOneTwentyK{} come within \NBudCloseGapOneTwentyK~pp of the mean wake loss of PSO-VNS (\NBudLossPSOVNSOneTwentyK\%). The ranking at 6,030 evaluations is thus that of a small budget.
% CHECK-FINAL [C26] [C29] [C34]: PSO-VNS best average rank at all three budgets; first-place counts generated per budget
% CHECK-FINAL [C37]: at 120,030 SSA-VNS ranks ahead of PSO; VNS, SSA-VNS, RS-VNS, LX-SSA-VNS and DE within 0.2 pp (\NBudCloseGapOneTwentyK) of PSO-VNS. PSO and VNS ranks are generated macros (not tested by C37).

\begin{figure}[!htbp]
\centering
\pipefig{0.8\textwidth}{figures_mpce/budget_scaling.pdf}{7cm}
\caption{Budget scaling: mean wake loss of the six largest benchmark cases and mean AEP loss of the Horns Rev~1 block (feasible runs; methods with at least half of their runs feasible) versus the budget (6,030, 30,030 and 120,030 evaluations, random initialization; 10 seeds for Horns Rev~1 at 120,030 evaluations).}
\label{fig:budget}
\end{figure}

Fig.~\ref{fig:budget} shows the case-level curves: the methods differ mainly in how much they gain from a larger budget. From 6,030 to 120,030 evaluations, the mean wake loss of VNS falls from \NBudLossVNSSixK\% to \NBudLossVNSOneTwentyK\%, that of PSO only from \NBudLossPSOSixK\% to \NBudLossPSOOneTwentyK\% (PSO-VNS: \NBudLossPSOVNSSixK\% to \NBudLossPSOVNSOneTwentyK\%), so the advantage of PSO-VNS over plain VNS is largest at the smallest budget. DE, feasible in only \NFbRandFeasDE\% of its runs at 6,030 evaluations, appears only at the larger budgets. On the Horns Rev~1 block (last panel), the mean AEP loss of PSO-VNS also falls with the budget (\NHRLossSixKPSOVNS\%, \NHRLossThirtyKPSOVNS\% and \NHRLossOneTwentyKPSOVNS\%; Table~\ref{tab:hr-site}), and at 120,030 evaluations \NHRBestHundredTwentyK{} has the lowest.
% verified (manual): budget losses from mpce_numbers.tex (\NBudLoss...): VNS - PSO-VNS mean wake loss 0.87 / 0.40 / 0.03 pp at 6,030 / 30,030 / 120,030 (decreasing); drop 6,030 -> 120,030: VNS 1.78, PSO 0.76, PSO-VNS 0.94 pp; DE feasible share at 6,030 from Table feasbudget (1.1 %); DE plotted only at 30,030 and 120,030 (at least half of the runs feasible)
% CHECK-FINAL [C25]: Horns Rev 16: 30,030 and 120,030 mean AEP loss of PSO-VNS below 6,030; CHECK-DIR [F29]: DE highest mean AEP (lowest loss) at 120,030; CHECK-FINAL [C37]: DE within 0.2 pp of PSO-VNS at 120,030

\subsection{Feasibility-Preserving Initialization}
\label{sec:feasinit}
With feasibility-preserving initialization, \NFeasInitPhrase. \NFbFeasBestRank{} alone is then best in mean wake loss and average rank, and PSO-VNS ranks second (rank \NFbFeasRankPSOVNS; loss \NFbFeasLossPSOVNS\%, \NFbRandLossPSOVNS\% with random starts; Table~\ref{tab:feasbudget}, per case: Table~\ref{S-tab:feasinit-cases}). PSO and DE never improve on the best initial layout (\NDFsStoredRunsPSO{} runs each): their moves change all or most coordinates at once and, between distinct dense layouts, make turbines collide (median minimum spacing after the first PSO move \NDFsFirstMinSpPSO~m, limit \NDSmin~m), so in our runs no trial is both feasible and better and the global best never moves (Lemma~\ref{lem:feas-improve}, Proposition~\ref{prop:feasible-start}; Table~\ref{tab:D-feasstart}). VNS, which moves one coordinate at a time, is not affected. The advantage of PSO-VNS thus depends on the initialization: from random starts, PSO may help mainly by reaching feasibility.
% CHECK-FINAL [C28] [C38]: feasible starts: VNS alone lowest loss and best rank; PSO-VNS second in rank; PSO and DE never move from the best initial layout
% CHECK-DIAG [D13]: PSO and DE never improve on the best initial layout (instrumented and all stored runs)
% CHECK-DIAG [D14]: all PSO candidates except the G re-evaluations and all \NDFsCandDE{} DE trials infeasible, none accepted, no pbest update; first-move median min spacing \NDFsFirstMinSpPSO{} m
% CHECK-DIAG [D16]: the feasible initial layouts are distinct (matched RMS distance to G > 50 m)

What follows from the algorithms can be separated from what is only observed. Let $L(\mathbf X)$ be the wake loss, the first term of~\eqref{eq:penalty}, so that $F_p=L+\mathrm{Pen}$ with $\mathrm{Pen}\ge0$.
\begin{lemma}[Improvement needs a feasible, better trial]
\label{lem:feas-improve}
If the global best $\mathbf G$ is feasible, a trial $\mathbf X$ replaces it only if $L(\mathbf X)<L(\mathbf G)$ and every constraint violation $g$ of $\mathbf X$ satisfies $(1+\mu g)^2<L(\mathbf G)$; in the benchmark cases, $g<4.6\cdot10^{-8}$ (m or m$^2$), far below the reporting tolerance of $10^{-6}$~m.
\end{lemma}
\begin{proof}
$\mathbf X$ replaces $\mathbf G$ only if $F_p(\mathbf X)<F_p(\mathbf G)=L(\mathbf G)$; as $L\ge0$ and each violation adds $(1+\mu g)^2>0$, each of these terms is below $L(\mathbf G)$. In the benchmark cases, $L(\mathbf G)$ is at most the wake-free objective $NE_{\rm ideal}\le15\times14{,}045.74$ plus a small term from the negative part of the power curve above cut-in, so $\sqrt{L(\mathbf G)}<459.1$ and, with $\mu=10^{10}$, $g<(\sqrt{L(\mathbf G)}-1)/\mu<4.6\cdot10^{-8}$.
\end{proof}
% verified: (sqrt(15 x 14045.74) - 1)/1e10 = 4.58e-8 (as optA/supp_theory.tex, Lemma S-feas-improve; the small negative part of the power curve above cut-in, epsilon, does not change the bound)
\begin{proposition}[PSO from feasible starts]
\label{prop:feasible-start}
Let all initial layouts be feasible, let the velocities start at zero with $\mathbf P^i=\mathbf X^i$, and let $i^*$ be the particle with the best initial layout, so that $\mathbf X^{i^*}=\mathbf P^{i^*}=\mathbf G$.
(a)~As long as $\mathbf G$ does not change, particle $i^*$ keeps $\mathbf V^{i^*}=\mathbf 0$ and $\mathbf X^{i^*}=\mathbf G$: it re-evaluates $\mathbf G$ once per iteration and evaluates no other point.
(b)~The first move of every other particle is $\mathbf X^i\leftarrow\mathbf X^i+c_2\boldsymbol\rho_2\circ(\mathbf G-\mathbf X^i)$, clipped to the box: each coordinate moves a random fraction in $[0,c_2]=[0,1.50]$ of the way to the same coordinate of $\mathbf G$, i.e., turbine $k$ of layout $i$ moves toward turbine $k$ of $\mathbf G$, although the numbering of the turbines of two independently generated layouts is arbitrary.
\end{proposition}
\begin{proof}
(a)~By induction: with $\mathbf V^{i^*}=\mathbf 0$ and $\mathbf P^{i^*}=\mathbf X^{i^*}=\mathbf G$, all three terms of~\eqref{eq:pso-v} vanish, the particle stays at $\mathbf G$, and $F_p(\mathbf G)$ is no strict improvement, so neither $\mathbf P^{i^*}$ nor $\mathbf G$ changes. (b)~At the first move, $\mathbf V^i=\mathbf 0$ and $\mathbf P^i=\mathbf X^i$, so only the last term of~\eqref{eq:pso-v} remains.
\end{proof}
DE is in a similar position: its trial takes coordinate $j$ from the mutant $\mathbf x_{r_1}+F(\mathbf x_{r_2}-\mathbf x_{r_3})$ if $j=j_{\rm rand}$ or $U_j<\mathrm{CR}$, so $1+\mathrm{Bin}(n-1,\mathrm{CR})$ coordinates, on average 27.1 of $n=30$ for $N=15$ ($\mathrm{CR}=0.9$), combine the same coordinates of three other layouts, again regardless of the turbine numbering. By Lemma~\ref{lem:feas-improve}, $\mathbf G$ moves only if such a combined layout is feasible and better; that this never happened in our runs is an observation, not a consequence of the proposition.
% CHECK-THEORY [T11]: DE takes most coordinates from the mutant (mean 1+CR(n-1) = 27.1 of 30 for N = 15; PSO moves all). Lemma lem:feas-improve = Lemma S-lem:S-feas-improve and Proposition prop:feasible-start = Proposition S-prop:S-frozen (a), (b) (optA/supp_theory.tex, S-th-feas); the stall is observed (D13/D14), not implied.
\SuppTable{tab:D-feasstart}

Table~\ref{tab:D-feasstart} logs every candidate (the six largest cases and the Horns Rev~1 block, \NDFsSeeds{} seeds each). Of the \NDFsCandPSO{} PSO candidates, \NDFsFeasCandPSO\% are feasible, and each of them is the particle holding $\mathbf G$ re-evaluating it, as in part~(a); no personal best is ever updated, and none of the \NDFsCandDE{} DE trials is feasible. Nearly all infeasible candidates violate both constraints (\NDFsInfBothPSO\% for PSO, \NDFsInfBothDE\% for DE). The arbitrary numbering of part~(b) is not the whole explanation: after relabelling the turbines of each start to their nearest turbines of $\mathbf G$, the first PSO move is feasible in only \NDFsMixMatched\% of the draws, and only a common weight for all coordinates (a convex combination of two layouts) reaches \NDFsMixScalar\%. The starts are distinct dense layouts (median RMS distance to $\mathbf G$ after relabelling \NDFsInitRms~m), and independent weights per coordinate, with overshoot beyond $\mathbf G$, place turbines too close together or outside the site.
% CHECK-DIAG [D14]: every feasible PSO candidate is the G-particle (zero velocity) re-evaluating G; no personal best updated; DE no feasible trial, none accepted
% CHECK-DIAG [D15]: after optimal relabelling the first PSO move is feasible in < 1 % of draws; only a common weight (convex combination) more often (\NDFsMixScalar %)
% CHECK-DIAG [D16]: feasible initial layouts distinct (matched RMS distance to G > 50 m; median \NDFsInitRms m)
% verified (manual): violation shares outside only / spacing only / both from Table D-feasstart (\NDFsInfBothPSO, \NDFsInfBothDE %)

\subsection{Horns Rev~1 16-Turbine Block}
\label{sec:hornsrev-site}
%% generated by mpce_results.py -- do not edit by hand
\begin{table}[!t]
\centering
\caption{Horns Rev~1 16-Turbine Block (Wake-Free AEP 149.26 GWh/yr): Mean AEP (GWh/yr) over the Feasible Runs of 30 Seeds at 6,030 Calls, Feasible Runs, Holm-Adjusted Wilcoxon $p$ of PSO-VNS vs.\ Each Method, and Mean AEP Loss (\%) at 6,030 Calls with Random (R) and Feasibility-Preserving (F) Initialization and at 30,030 and 120,030 Calls (R; 10 Seeds at 120,030); Superscript: Feasible Runs When Not All Runs Are Feasible}
\label{tab:hr-site}
\scriptsize\setlength{\tabcolsep}{2.2pt}
\begin{tabular}{lccccccc}
\toprule
& \multicolumn{3}{c}{6,030 calls, R} & \multicolumn{4}{c}{Loss (\%)} \\
\cmidrule(lr){2-4}\cmidrule(lr){5-8}
Layout / method & AEP & Feas. & $p_{\rm Holm}$ & 6k R & 6k F & 30k & 120k \\
\midrule
Installed layout & 139.51 & -- & -- & 6.53 & -- & -- & -- \\
\midrule
PSO-VNS & 138.73 & 30/30 & -- & 7.05 & 6.96 & 6.46 & \TBD{} \\
SSA-VNS & 137.98 & 30/30 & $6.7\times10^{-4}$ & 7.56 & \TBD{} & \TBD{} & \TBD{} \\
LX-SSA-VNS & 137.77 & 30/30 & $6.3\times10^{-4}$ & 7.70 & \TBD{} & \TBD{} & \TBD{} \\
RS-VNS & 137.96 & 28/30 & $2.2\times10^{-4}$ & 7.57$^{28}$ & \TBD{} & \TBD{} & \TBD{} \\
VNS & 137.67 & 30/30 & $2.2\times10^{-4}$ & 7.77 & \TBD{} & \TBD{} & \TBD{} \\
PSO & 138.06 & 19/30 & $1.4\times10^{-5}$ & 7.51$^{19}$ & \TBD{} & \TBD{} & \TBD{} \\
SSA & 136.66 & 30/30 & $2.8\times10^{-7}$ & 8.44 & \TBD{} & \TBD{} & \TBD{} \\
LX-SSA & 136.46 & 28/30 & $1.5\times10^{-7}$ & 8.58$^{28}$ & \TBD{} & \TBD{} & \TBD{} \\
DE & -- & 0/30 & $1.7\times10^{-8}$ & -- & \TBD{} & \TBD{} & \TBD{} \\
MS-SLSQP & 136.80 & 30/30 & $1.5\times10^{-7}$ & 8.35 & \TBD{} & \TBD{} & \TBD{} \\
\bottomrule
\end{tabular}
\end{table}


The model (Section~\ref{S-sec:S-hrmodel}) uses the Jensen wake with a hub-center in-wake test, the thrust coefficient $C_T$ at the free-stream speed, and 5\textdegree{} direction and 1-m/s speed bins; the boundary violation in~\eqref{eq:penalty} is the distance (m) outside the parallelogram. For the installed 80-turbine farm at identical bins (same wake-free AEP), PyWake~\NFPyWakeVersion's NOJ model~\cite{PyWake} gives \NFPyWakeAEP~GWh/yr and ours \NFPyWakeOurAEP~GWh/yr (\NFPyWakeDiffPct\% more; wake loss \NFPyWakeOurLossPct\% vs \NFPyWakeLossPct\%; Table~\ref{S-tab:F-pywake}). Rotor averaging does not explain the difference (PyWake's hub-center variant lies further from our value); with $C_T$ at the local speed, our model reproduces PyWake's hub-center NOJ with momentum-theory induction exactly. Our model thus slightly underestimates the wake loss (installed 16-turbine block: AEP \NFPyWakeSixteenDiffPct\% higher); the same model is used for all methods.
% CHECK-DIR [F31] [F32]: PyWake 2.6.20 NOJ(k=0.04) at bins identical to ours (same wake-free AEP): 661.39 vs ours 666.75 GWh/yr (+0.81 %), wake loss 11.11 vs 10.39 %
% CHECK-DIR [F33] [F34]: rotor averaging does not explain the difference; with C_T at the local (waked) speed our model reproduces PyWake's NOJ with hub-centre test (momentum-theory induction) to 1e-6 GWh/yr; the old 662.5 constant (\NHRPyWakeDiff) is dropped [F36]
% CHECK-DIR [F35]: installed 16-turbine block: PyWake NOJ 139.08 vs ours 139.82 GWh/yr (+0.53 %)

The installed layout, a regular $7D$ grid, reaches \NHRInstalled~GWh/yr (wake loss \NHRInstalledLoss\%; Table~\ref{tab:hr-site}). At 6,030 evaluations from random starts, the robust result is feasibility: PSO-VNS is feasible in \NHRFeasPSOVNS{} runs, PSO in only \NHRFeasPSO{} (exact McNemar $p=\NXHRMcNemarP$; \NHRFeasFeasInitPSO{} with feasible starts) and DE in \NHRFeasDE. The best layout of all methods at this budget, a PSO-VNS layout (\NHRBestPSOVNS~GWh/yr; Fig.~\ref{fig:layouts-sites}, right), stays below the installed block, like every run at this budget; the convergence curves are in Fig.~\ref{S-fig:S-hr16}.
% CHECK-EXTRA [X52]: HR feasibility PSO-VNS 30/30 vs PSO 19/30, 11 discordant seeds all favour PSO-VNS, exact McNemar p = 9.8e-4
% CHECK-FINAL [C20]: PSO feasible in fewer than 30 of 30 random-start runs ("only"); [C31]: 6,030 PSO-VNS mean AEP below the installed block; no 6,030 random-start run above the installed block with 5-deg bins (Table F-hr; 16 runs above at 5 deg over all budgets, 0 at 1 deg: [F21])
% Over the page budget (+0.35 pages), kept in the supplement (Fig. S-fig:S-hr16):
% \begin{figure}[!htbp]
% \centering
% \pipefig{0.75\textwidth}{figures_mpce/hr16.pdf}{6cm}
% \caption{Horns Rev~1 16-turbine block (5\textdegree{} bins, 6,030 evaluations, random initialization): median best feasible AEP versus evaluations (left) and the best layout of all methods (right).}
% \label{fig:hr16}
% \end{figure}

We also re-evaluated every final layout with 1\textdegree{} bins and with PyWake's NOJ at the same bins (Table~\ref{S-tab:F-hr}). The AEP advantage of PSO-VNS over PSO at 6,030 evaluations on the \NFHRPairSixKOneN{} jointly feasible seeds (\NFHRPairSixKFive~GWh/yr, $p=\NFHRPairSixKFiveP$) vanishes with 1\textdegree{} bins (\NFHRPairSixKOne, $p=\NFHRPairSixKOneP$) and under PyWake. At 30,030 evaluations, PSO-VNS has the highest mean AEP under all three evaluations (5\textdegree{} bins: \NHRMeanThirtyKPSOVNS~GWh/yr) and beats PSO on the \NFHRPairThirtyKOneN{} jointly feasible seeds also with 1\textdegree{} bins (\NFHRPairThirtyKOne~GWh/yr, $p=\NFHRPairThirtyKOneP$) and under PyWake; at 120,030 evaluations (10 seeds), \NHRBestHundredTwentyK{} has the highest mean AEP. The finer rose costs the installed block \NFHRInstDrop~GWh/yr and the optimized layouts \NFHRMeanDrop~GWh/yr on average, the VNS-based methods most: like those of the benchmark (Section~\ref{sec:direction}), the optimized layouts exploit the gaps between the coarse direction bins. With 1\textdegree{} bins, \NFHRAboveOneTotal{} of the \NFHRNFeas{} feasible optimized layouts (all methods and budgets) exceed the installed block (\NFHRAboveFiveTotal{} with 5\textdegree{} bins; under PyWake \NFHRAbovePyWakeTotal, a DE run at 120,030), although the looser spacing $4D$ (installed: $7D$) enlarges their feasible set. The block is evaluated without the other 64 turbines, and cables, foundations, loads and navigation are not modeled: the comparison tests the optimizers, not the installed design.
% CHECK-DIR [F24]: 6,030 PSO-VNS - PSO AEP on 19 jointly feasible seeds: significant only with 5-deg bins; with 1-deg bins and PyWake NOJ PSO has the (n.s.) higher mean
% CHECK-DIR [F26]: 30,030: PSO-VNS highest mean AEP with 5-deg, 1-deg (both phases) and PyWake NOJ; beats PSO on 28 jointly feasible seeds with 1-deg bins and PyWake (p 6.5e-4, 3.8e-4)
% CHECK-DIR [F29]: 120,030: DE highest mean AEP under 5-deg bins, 1-deg bins and PyWake NOJ
% CHECK-DIR [F23]: installed block loses \NFHRInstDrop GWh/yr, optimized layouts \NFHRMeanDrop on average; every VNS-based method loses more than every other method
% CHECK-DIR [F21] [F22]: with 1-deg bins 0 of 901 feasible runs exceed the installed block (16 with 5-deg bins); PyWake NOJ: 1 run (DE, 120,030). The former "runs exceed the installed layout" (C30) and the 6,030 run-level test / highest mean AEP (C19, C39) are no longer asserted (R4-1, R4-10; D16, D19).
% Over the page budget (+1.0 and +0.6 pages in the private build), kept in the supplement: \SuppTable{tab:F-hr} \SuppTable{tab:F-pywake}

\subsection{IEA37 Case Study~1}
\label{sec:iea37}
%% generated by mpce_results.py -- do not edit by hand
\begin{table}[!t]
\centering
\caption{IEA37 Case Study~1, 16- and 36-Turbine Scenarios: AEP in MWh of the Best Run of Each Method at 6,030 and 30,030 Calls (30 Seeds), Compared with the Example Layout and the Best Feasible Published Layout~\cite{Baker2019,IEA37repo}}
\label{tab:iea37}
\scriptsize\setlength{\tabcolsep}{2.5pt}
\begin{tabular}{lcccc}
\toprule
& \multicolumn{2}{c}{16 turbines ($r=1300$~m)} & \multicolumn{2}{c}{36 turbines ($r=2000$~m)} \\
\cmidrule(lr){2-3}\cmidrule(lr){4-5}
Method / source & 6,030 & 30,030 & 6,030 & 30,030 \\
\midrule
Example layout & \multicolumn{2}{c}{366{,}941.6} & \multicolumn{2}{c}{737{,}883.1} \\
Best feasible published$^{a}$ & \multicolumn{2}{c}{418{,}924.4} & \multicolumn{2}{c}{863{,}676.3} \\
\midrule
PSO-VNS & \TBD{} & \TBD{} & \TBD{} & \TBD{} \\
PSO & 403{,}753.1 & 412{,}702.4 & 787{,}988.7 & 830{,}256.4 \\
SSA-VNS & 403{,}251.3 & 406{,}924.6 & 766{,}803.3 & 819{,}930.5 \\
SSA & 395{,}150.2 & 401{,}459.3 & 750{,}707.4 & 802{,}538.2 \\
LX-SSA & 393{,}284.6 & 398{,}770.0 & 746{,}301.3 & 789{,}311.3 \\
DE & 357{,}391.3 & 397{,}402.7 & -- & -- \\
VNS & 400{,}169.3 & 404{,}076.0 & 785{,}791.3 & 824{,}230.2 \\
MS-SLSQP & 405{,}433.8 & 410{,}349.5 & 748{,}571.1 & 833{,}365.0 \\
\bottomrule
\multicolumn{5}{p{0.92\columnwidth}}{$^{a}$Participant~4 in both scenarios (SNOPT with wake expansion continuation \TBD{verify against Baker et al.\ (2019)}). A higher submitted AEP places turbines up to 3.5~m outside the boundary and is not counted.}
\end{tabular}
\end{table}


In IEA37 Case Study~1 the participants chose their own optimizers and budgets~\cite{Baker2019}. None of our methods reaches the best feasible published layouts (Table~\ref{tab:iea37}; Fig.~\ref{fig:layouts-sites}). Table~\ref{tab:iea37} counts a published layout as feasible at a tolerance of \NFIEATolStrict{} (ours: $10^{-6}$~m); after radial projection onto the boundary, \NFIEANChanged{} more are feasible, and the best are then those of participant~12. Against these, the best PSO-VNS layout at 30,030 evaluations differs in AEP by \NFIEAGapSixteenThirtyKProj\% (16 turbines) and \NFIEAGapThirtySixThirtyKProj\% (36 turbines; strict: \NFIEAGapSixteenThirtyKStrict\% and \NFIEAGapThirtySixThirtyKStrict\%), in wake loss by \NFIEAGapLossSixteenThirtyKProj{} and \NFIEAGapLossThirtySixThirtyKProj~pp, and would rank \NFIEARankSixteenThirtyKProj{} among \NFIEAPubFeasProjSixteen{} and \NFIEARankThirtySixThirtyKProj{} among \NFIEAPubFeasProjThirtySix{} feasible submissions. With 36 turbines, the best of our layouts comes from \NFIEAOurBestByThirtySixThirtyK{} (\NFIEAOurBestGapThirtySixThirtyKProj\%). As the participants' budgets are not reported comparably, the size of the gap is indicative only. The 64-turbine scenario was not run, and, like the published layouts, ours are optimized for 16 discrete directions and one wind speed.
% CHECK-DIR [F41] [F42]: 5 published layouts violate only the boundary (all but one by < 1 cm) and are feasible after radial projection; the best feasible published layout is then participant 12 in both scenarios (421,451.1 and 882,382.8 MWh) instead of participant 4
% CHECK-DIR [F43] [F44]: strict gaps reproduce \NIEA... (-1.33 / -4.20 %); projected gaps of the best PSO-VNS layout at 30,030: -1.93 / -6.23 % (wake loss +1.73 / +5.21 pp); rank fourth of 12+1 and ninth of 10+1
% CHECK-DIR [F45] + CHECK-FINAL [C48]: 36 turbines, 30,030: best of our layouts from MS-SLSQP (-5.56 % projected)
% CHECK-FINAL [C40]: best of all methods below the best feasible published layout (strict) in both scenarios

Table~\ref{tab:F-iea-gap} gives all gaps. At 6,030 evaluations they are larger (best PSO-VNS layout, projected: \NFIEAGapSixteenSixKProj\% and \NFIEAGapThirtySixSixKProj\%), so the gap depends more on the budget than on the feasibility convention, and it is larger for 36 than for 16 turbines at both budgets and under both conventions.
% CHECK-DIR [F43] [F44]: gaps of Table F-iea-gap (generated macros \NFIEAGap..., \NFIEAMeanGap...); CHECK-FINAL [C40]: gap larger for 36 turbines
% verified (manual): budget effect on the best PSO-VNS gap (projected) 2.50 / 4.50 pp vs convention effect (30,030) 0.60 / 2.03 pp
\SuppTable{tab:F-iea-gap}
\begin{figure}[!htbp]
\centering
\pipefig{\textwidth}{figures_mpce/layouts_iea37_hr16.pdf}{5cm}
\caption{Best layouts. Left and middle: IEA37 Case Study~1, 16 and 36 turbines, best feasible published layout (strict convention) and best PSO-VNS layout at 30,030 evaluations. Right: Horns Rev~1 16-turbine block (5\textdegree{} bins), installed layout ($7D$ grid) and best PSO-VNS layout at 6,030 evaluations (minimum spacing $4D$).}
\label{fig:layouts-sites}
\end{figure}

\subsection{Computational Cost and Published Results}
\label{sec:literature}
\label{sec:runtime}
Per evaluation, the metaheuristics need \NMsEvalMin--\NMsEvalMax~ms and MS-SLSQP \NMsEvalSLSQP~ms, about \NSLSQPSlowdown{} times as long as PSO (indicative; Table~\ref{S-tab:cost}); one 6,030-evaluation run of PSO-VNS takes \NSecPerRunMin--\NSecPerRunMax~s, depending on $N$. Published Kusiak--Song results~\cite{Kusiak2010,Eroglu2012,Eroglu2013,Bansal2017} are discussed in Section~\ref{S-sec:S-literature}; they are not compared numerically because budgets, constraint handling and reporting differ.
% CHECK-FINAL [C46]: cost per evaluation (metaheuristics similar, MS-SLSQP slower than every metaheuristic)
